\documentclass{ximera}

% MATH -----------------------------------------------------------
\newcommand{\nats}{\mathbb N}
\newcommand{\ints}{\mathbb Z}
\newcommand{\rats}{\mathbb Q}
\newcommand{\reals}{\mathbb R}
\newcommand{\complex}{\mathbb C}
\newcommand{\powerset}{\mathscr P}

%-------------------------------------------------------------

\pgfplotsset{compat=1.18}
\begin{document}
\begin{abstract}
 \end{abstract}

    \title{Exemplar Examples 18}
    
    \maketitle


Suppose a mass is hung on a large vertical spring and is brought to rest at equilibrium, causing the spring to stretch from its natural length (with no mass attached) by 2.45 centimeters.

\begin{enumerate}

\item Let's determine the natural angular frequency of this spring mass system.  What is the frequency it in hertz?  We will use $g=9.8$ m$/s^2$ for the magnitude of the acceleration due to gravity.

\vfill


% w=sqrt(g/l)=sqrt(980/2.45)=20 rad/s ----> f=w/(2pi)=10/pi hz (cycles/s)


\item Suppose the mass is pulled down by 5 centimeters and released.  Assuming it is free and undamped, let's find the displacement $y(t)$ of the mass above equilibrium (in centimeters) at $t\geq 0$ seconds after being released.  Let's use this to determine the velocity of the mass at $t=\dfrac{\pi}{40}$ seconds (as the mass passes through the equilibrium position for the first time).

\vfill
\vfill

% y=-5\cos(20t).  y'=100 sin(20t)  ... 100 cm/s [OR 1 m/s]


\newpage

\item Now let's assume the mass is $1$ kg and that there is a mechanical damper attached to the spring which damps the motion at a rate of $40$ N per m/s of speed.  Using the same conditions (the mass is pulled down by 5 centimeters and released) let's determine the displacement $y(t)$ of the mass above equilibrium (in centimeters) at $t\geq 0$ seconds after being released.  Is this system over-damped, under-damped, or critically damped.  Let's determine the number of times the mass passes through equilibrium.

\vfill

\item Now let's assume that in addition to the mechanical damper, we are forcing the mass with a constant $20$ Newtons of upward force.  Using the same conditions (the mass is pulled down by 5 centimeters and released) let's determine the displacement $y(t)$ of the mass above equilibrium (in centimeters) at $t\geq 0$ seconds after being released.

\vfill


\end{enumerate}

\end{document}